\section*{Chapter \ref{ch:b1:radiation-heat-transport} --- Radiation and Heat Transport}
% values: calc/ch27.py

\begin{solution}{exo:b1:radiation-heat-transport:1}
(a)~$S_0/\pi = 1361/\pi = \qty{433.2}{W.m^{-2}}$. (b)~At the equator on an equinox
$\bar Q = (S_0/\pi)(a/r)^2$, so $(a/r)^2 = 436.7/433.2 = 1.0080$ and $r/a =
1.0080^{-1/2} = 0.9960$: the Earth is \qty{0.4}{\percent} closer to the Sun than
its mean distance. (c)~$\bar Q = 433.2 \times 0.9683 \times \cos\ang{23.44} =
\qty{384.9}{W.m^{-2}}$.
\end{solution}

\begin{solution}{exo:b1:radiation-heat-transport:2}
(a)~$415.6/172.4 = 2.41$. (b)~The zone between the parallels $\pm\varphi$ has the
area $2\pi a_\oplus^2 \times 2\sin\varphi$, a fraction $\sin\varphi$ of the sphere:
$\sin\ang{37.0} = 0.602$, \qty{60}{\percent}. (c)~On an axis in \omterm{def:b1:spherical-astronomy:geographic}{latitude}, the
bands near the poles look as wide as those near the equator, but they cover much
less area: the zone poleward of \ang{60} is only $1 - \sin\ang{60} =
\qty{13.4}{\percent}$ of the surface. Against $\sin\varphi$, the tropics fill most
of the axis.
\end{solution}

\begin{solution}{exo:b1:radiation-heat-transport:3}
(a)~$T_e = (\text{OLR}/\sigma)^{1/4}$: \qty{256.4}{K}, \qty{241.3}{K} and
\qty{254.5}{K}. (b)~Absorbed sunlight: $287.7/75.8 = 3.80$; OLR: $244.9/192.3 =
1.27$. (c)~The \omterm{def:b1:radiation-heat-transport:net-radiation}{net radiation} follows the absorbed sunlight: positive and large in the
tropics, strongly negative at high \omterm{def:b1:spherical-astronomy:geographic}{latitudes}. The OLR, which varies little, does
not offset the fall of the absorbed sunlight.
\end{solution}

\begin{solution}{exo:b1:radiation-heat-transport:4}
(a)~$2.82 \times 365.25 = \qty{1030}{mm}$, \qty{1.03}{m} a year. (b)~$1.030 \times
78.1 = \qty{80.4}{W.m^{-2}}$.
(c)~The reanalysis's $LE$ is \qty{2}{W.m^{-2}} larger. Over many years the
atmosphere cannot keep accumulating water vapour: all the water evaporated comes
down again, so $LE$ and $LP$ must be equal in the global mean. The difference is
an imbalance of the model, like its radiative imbalance.
\end{solution}

\begin{solution}{exo:b1:radiation-heat-transport:5}
(a)~$\bar Q = S_0(a/r)^2\sin\delta = 1361 \times 0.9683 \times 0.3978 =
\qty{524.2}{W.m^{-2}}$. (b)~$1361 \times 1.0334 \times 0.3978 =
\qty{559.5}{W.m^{-2}}$; the ratio is $1.0334/0.9683 = 1.067$. (c)~The pole has
sunlight 24 hours out of 24 at a constant \omterm{def:b1:spherical-astronomy:horizontal}{altitude}, so its mean $\sin h_\odot$ over
the day is $\sin\delta = 0.398$. At the equator the Sun is high at noon but down
half the time; the mean of $\sin h_\odot$ over the 24 hours is $\cos\delta/\pi =
0.292$. The ratio is $\pi\tan\delta = 1.362$
(\cref{prop:b1:radiation-heat-transport:polar-summer}).
\end{solution}

\begin{solution}{exo:b1:radiation-heat-transport:6}
(a)~$\cos h_0 = -\tan\ang{40}\tan\ang{23.44} = -0.3638$, $h_0 = \ang{111.33} =
\qty{1.943}{rad}$; the day lasts $24 \times 111.33/180 = \qty{14.84}{h}$; $\bar Q =
(1361/\pi) \times 0.9683 \times (1.943 \times 0.6428 \times 0.3978 + 0.7660 \times
0.9175 \times 0.9315) = \qty{483.0}{W.m^{-2}}$. (b)~$\cos h_0 = +0.3638$, $h_0 =
\ang{68.67}$, \qty{9.16}{h}; $\bar Q = \qty{155.9}{W.m^{-2}}$. (c)~The ratio is
$483.0/155.9 = 3.10$. The distance contributes $0.9683/1.0334 = 0.937$: it works
against the northern summer. The geometry, the \omterm{def:b1:spherical-astronomy:horizontal}{altitude} of the Sun and the length
of the day, contributes $3.10/0.937 = 3.31$, all of the seasonal contrast and more.
\end{solution}

\begin{solution}{exo:b1:radiation-heat-transport:7}
(a)~$B = H/LE$: \num{0.093} at the equator, \num{0.374} at \ang{25} N,
\num{0.338} at \ang{45} N, \num{-0.152} at \ang{70} N. (b)~\ang{25} N: the zonal
mean includes the great hot deserts, dry land with little water to evaporate, where
the heated ground warms the air instead. (c)~$H < 0$ means that \omterm{def:b1:heat-budget:heat-flow}{heat flows} from the
air to the surface: in the annual mean the snow and sea ice are colder than the air
above them, which is warmed by air brought from lower \omterm{def:b1:spherical-astronomy:geographic}{latitudes}.
\end{solution}

\begin{solution}{exo:b1:radiation-heat-transport:8}
(a)~Assessed: $\alpha = 100/340 = 0.294$, $R = 340 - 100 - 239 =
\qty{+1}{W.m^{-2}}$ (the \num{0.7} rounded). Reanalysis: $\alpha = 117.0/341.9 =
0.342$, $R = 341.9 - 117.0 - 237.7 = \qty{-12.8}{W.m^{-2}}$ with these rounded
terms, \qty{-12.9}{W.m^{-2}} with the unrounded ones of
\cref{ex:b1:radiation-heat-transport:global}. (b)~The reanalysis
absorbs $224.8$ against $240$, \qty{15.2}{W.m^{-2}} less, because it reflects
\qty{17}{W.m^{-2}} more; its OLR is only \qty{1.3}{W.m^{-2}} lower. The excess
reflection explains the imbalance. (c)~Integrating $R$ from the South Pole to the
North Pole gives the global integral, $-12.9 \times 4\pi a_\oplus^2 = -12.9 \times
5.10\times10^{14} = \qty{-6.6}{PW}$. A transport across the North Pole is
impossible: there is nothing beyond it, the parallel is a point. The proposition
holds only if $\int R\,\mathrm dA = 0$, which the steady state requires.
\end{solution}

\begin{solution}{exo:b1:radiation-heat-transport:9}
(a)~The area is $2\pi a_\oplus^2 \times 2\sin\ang{30} = 2\pi a_\oplus^2 =
\qty{2.55e14}{m^2}$, half the Earth; the surplus is $33 \times 2.55\times10^{14} =
\qty{8.4}{PW}$. (b)~\qty{4.2}{PW} across each parallel. (c)~The two add up to
\qty{8.5}{PW}, the surplus, but they are unequal: the southern hemisphere exports
\qty{0.5}{PW} more than the northern one. The surplus is not shared at the
equator: in this reanalysis \qty{0.36}{PW} of the northern tropics' surplus crosses
the equator southwards.
\end{solution}

\begin{solution}{exo:b1:radiation-heat-transport:10}
(a)~Write the daily mean as $\bar Q = (a/r)^2 F(\varphi, \delta)$, with $F$ the
bracket of \cref{thm:b1:radiation-heat-transport:daily} times $S_0/\pi$. During
$\mathrm dt$ the energy received per square metre is $\bar Q\,\mathrm dt$, and
$\mathrm dt = r^2\,\mathrm d\lambda/h$, so $\bar Q\,\mathrm dt = (a^2/h)
F(\varphi, \delta(\lambda))\,\mathrm d\lambda$: the factor $r^2$ cancels.
(b)~Over a year the energy is $(a^2/h)\int_0^{2\pi} F(\varphi,
\delta(\lambda))\,\mathrm d\lambda$. The formula is unchanged when $\varphi$ and
$\delta$ change sign together, $F(-\varphi, \delta) = F(\varphi, -\delta)$, and
$\delta(\lambda + \pi) = -\delta(\lambda)$, since $\sin\delta = \sin\varepsilon
\sin\lambda$. Shifting $\lambda$ by $\pi$ in the integral shows that $-\varphi$
receives the same energy as $\varphi$. (c)~The northern spring and summer last
\num{7.6} days longer, because the Earth is near aphelion and moves slowly then; the
southern summer is brighter but shorter. By (a), the two effects cancel exactly.
\end{solution}

\begin{solution}{exo:b1:radiation-heat-transport:11}
(a)~$\langle (a/r)^2\rangle = \frac1P\int_0^P \frac{a^2}{r^2}\,\mathrm dt =
\frac1P\int_0^{2\pi}\frac{a^2}{r^2}\frac{r^2}{h}\,\mathrm d\lambda = \frac{2\pi
a^2}{hP} = \frac{2\pi a^2}{2\pi a b} = \frac ab$, and $b = a\sqrt{1-e^2}$.
(b)~The instantaneous global mean $\frac14 S_0(a/r)^2$ averages to $\frac14 S_0
\langle(a/r)^2\rangle = S_0/(4\sqrt{1-e^2})$. (c)~$340.25/\sqrt{1 - 0.0167^2} =
\qty{340.30}{W.m^{-2}}$, \qty{0.05}{W.m^{-2}} or \qty{0.014}{\percent} above
$S_0/4$.
\end{solution}

\begin{solution}{exo:b1:radiation-heat-transport:12}
(a)~Atmosphere $0.78 \times 4.10 = \qty{3.20}{PW}$, ocean $0.22 \times 4.10 =
\qty{0.90}{PW}$. (b)~The atmosphere alone is found to carry more, \qty{5.0}{PW},
than the reanalysis's atmosphere and ocean together, \qty{4.02}{PW}: the
reanalysis's total transport is short by at least a fifth ($1 - 4.02/5.0 = 0.20$).
Its \omterm{def:b1:radiation-heat-transport:net-radiation}{net radiation} varies too little between the tropics and the poles. (c)~Removing a
constant from $R$ makes its global integral zero but leaves its variation with
\omterm{def:b1:spherical-astronomy:geographic}{latitude} unchanged; the error lies in that variation, in the albedo of the model's
clouds and surfaces at each \omterm{def:b1:spherical-astronomy:geographic}{latitude}, not only in the global mean.
\end{solution}

\begin{solution}{pb:b1:radiation-heat-transport:1}
\begin{enumerate}
  \item June: $\cos h_0 = -0.3638$, $h_0 = \ang{111.33}$, $\bar Q =
        \qty{483.0}{W.m^{-2}}$. December: $h_0 = \ang{68.67}$, $\bar Q =
        \qty{155.9}{W.m^{-2}}$. March: $h_0 = \ang{90}$, $\bar Q = (1361/\pi)
        \times 1.0081 \times \cos\ang{40} = \qty{334.6}{W.m^{-2}}$.
  \item June: $\tan\ang{80}\tan\ang{23.44} = 2.46 > 1$, polar day, $h_0 = \pi$:
        $\bar Q = 1361 \times 0.9683 \times \sin\ang{80}\sin\ang{23.44} =
        \qty{516.3}{W.m^{-2}}$. December: polar night, $\bar Q = 0$. March:
        $(1361/\pi) \times 1.0081 \times \cos\ang{80} = \qty{75.8}{W.m^{-2}}$.
  \item At \ang{80} N the Sun does not set ($h_0 = \ang{180}$, 24 hours) while at
        \ang{40} N it shines for \qty{14.8}{h}. Its \omterm{def:b1:spherical-astronomy:horizontal}{altitude} is lower, between
        \ang{13.4} and \ang{33.4}, but the term $h_0\sin\varphi\sin\delta$, which
        grows with the day length and the \omterm{def:b1:spherical-astronomy:geographic}{latitude}, wins.
  \item The band values exceed the central values by \qty{0.5}{\percent},
        \qty{0.8}{\percent}, \qty{1.7}{\percent}, \qty{1.7}{\percent} and
        \qty{1.3}{\percent}. First, the reanalysis uses a slightly larger solar
        irradiance (Question~5), \qty{0.5}{\percent} more. Second, a band mean is
        not the value at the centre: the equatorward half of each band has more
        area and more sunlight than the poleward half.
  \item By \cref{prop:b1:radiation-heat-transport:global}, $S_0 = 4\sqrt{1-e^2}
        \times 341.9 = \qty{1367}{W.m^{-2}}$, \qty{0.47}{\percent} more than \qty{1361}{W.m^{-2}}.
  \item $A = 2\pi a_\oplus^2(1 - \sin\ang{40}) = 2\pi \times (6.371\times10^6)^2
        \times 0.3572 = \qty{9.110e13}{m^2}$, a fraction $(1 - \sin\ang{40})/2 =
        0.179$ of the Earth.
  \item $(1 - \alpha)\bar Q$: \num{198.0}, \num{159.3}, \num{124.7}, \num{92.5},
        \num{74.7} \unit{W.m^{-2}}.
  \item $R$: \num{-28.1}, \num{-51.0}, \num{-76.6}, \num{-104.9}, \num{-116.9}
        \unit{W.m^{-2}}. All five bands are in deficit; the polar band
        \ang{80}--\ang{90} has the largest.
  \item $T_e = (\text{OLR}/\sigma)^{1/4}$: \num{251.3}, \num{246.8}, \num{244.1},
        \num{242.9}, \num{241.1} \unit{K}.
  \item Absorbed: $198.0/74.7 = 2.65$; OLR: $226.1/191.6 = 1.18$.
  \item The Sun is lower, and light arriving at a grazing angle is reflected more
        by clouds and water; snow and sea ice, bright surfaces, cover more of the
        ground and the sea towards the pole.
  \item The OLR goes as $T_e^4$, and $T_e$ varies by only \qty{10}{K}, \qty{4}{\percent},
        across the cap: the high \omterm{def:b1:spherical-astronomy:geographic}{latitudes} are kept much warmer than their sunlight
        alone would allow, by the heat brought from the south.
  \item In \unit{m^2}: \num{3.143e13}, \num{2.550e13}, \num{1.879e13},
        \num{1.151e13}, \num{0.387e13}; the sum is \qty{9.110e13}{m^2}.
  \item $R_iA_i$: \num{-0.883}, \num{-1.300}, \num{-1.439}, \num{-1.207},
        \num{-0.453} \unit{PW}; the cap: \qty{-5.28}{PW}.
  \item $\alpha_\text{cap} = \sum\alpha_i\bar Q_iA_i/\sum\bar Q_iA_i = 0.400$,
        against \num{0.342} for the globe.
  \item $-12.90 \times 5.10\times10^{14} = \qty{-6.58}{PW}$: the Earth would be
        losing heat at that rate and cooling. The assessed imbalance is
        $+0.7 \times 5.10\times10^{14} = \qty{+0.36}{PW}$, eighteen times smaller and
        of the opposite sign: the reanalysis's imbalance is a model error.
  \item Corrected $R_i + 12.90$: \num{-15.2}, \num{-38.1}, \num{-63.7},
        \num{-92.0}, \num{-104.0} \unit{W.m^{-2}}; powers \num{-0.478},
        \num{-0.972}, \num{-1.197}, \num{-1.059}, \num{-0.403} \unit{PW}; the cap:
        \qty{-4.108}{PW}, a mean deficit of \qty{45.1}{W.m^{-2}}.
  \item The cap's energy does not change from year to year, so what it loses at
        the top of the atmosphere must enter it, and the only way in is across the
        40th parallel: $T(\ang{40}) = -2\pi a_\oplus^2\int_{\ang{40}}^{\ang{90}}
        R\cos\varphi\,\mathrm d\varphi$.
  \item $T(\ang{40}\,\text{N}) = +\qty{4.11}{PW}$, northwards.
  \item $LP = 2.4655\times10^6 \times 1.95/86400 = \qty{55.6}{W.m^{-2}}$, more than
        $LE = \qty{42.4}{W.m^{-2}}$: more water falls on the cap than evaporates
        from it; the difference arrives as vapour.
  \item $(LE - LP) - 2.07 = 42.4 - 55.6 - 2.07 = \qty{-15.3}{W.m^{-2}}$; times the
        area of the cap: \qty{-1.40}{PW}. The atmosphere brings \qty{1.40}{PW}
        of latent heat across \ang{40} N.
  \item $1.40/4.11 = 0.34$: a third.
  \item Ocean $0.22 \times 4.11 = \qty{0.90}{PW}$; atmosphere $0.78 \times 4.11 =
        \qty{3.20}{PW}$, of which \qty{1.40}{PW} latent and \qty{1.80}{PW}
        carried as sensible heat and potential energy of the air.
  \item A lower bound. The reanalysis's total at \ang{43} N, \qty{4.02}{PW}, is less
        than the atmosphere alone in the study, \qty{5.0}{PW}; the true total is
        larger.
  \item About \qty{4.11}{PW}, \qty{4.11e15}{W}, according to the reanalysis's
        radiation, a third of it as latent heat; the true value is larger still.
\end{enumerate}
\end{solution}
